% ============================================================================
% DeepCatch: Performance-Weighted Multi-Modal Fusion with Cumulative Evidence
% Tracking for Ultra-Early Cancer Detection from Cell-Free DNA
%
% Target Journal: Bioinformatics (Oxford Academic) — Computational Methods
% Alternate: PLOS Computational Biology
% Type: Original Article (Methods)
% Date: April 28, 2026
% ============================================================================

%\documentclass{bioinfo}
\documentclass[10pt,a4paper]{article}

\usepackage[utf8]{inputenc}
\usepackage[T1]{fontenc}
\usepackage{amsmath,amssymb,amsfonts}
\usepackage{graphicx}
\usepackage{booktabs}
\usepackage{hyperref}
\usepackage{natbib}
\usepackage[margin=2.5cm]{geometry}
\usepackage{lineno}
\usepackage{xcolor}
\usepackage{float}
\usepackage{caption}
\usepackage{subcaption}
\usepackage{multirow}
\usepackage{array}
\usepackage{enumitem}

\linenumbers

\title{\textbf{DeepCatch: Performance-Weighted Multi-Modal Fusion with Cumulative Evidence Tracking for Ultra-Early Cancer Detection from Cell-Free DNA}}

\author{
  Royce [Surname]$^{1,*}$ and the DeepCatch Contributors\\
  \small{$^{1}$Independent Researcher}\\
  \small{$^{*}$Corresponding author: \texttt{royce@deepcatch.org}}
}

\date{}

\begin{document}

\maketitle

% ============================================================================
% ABSTRACT
% ============================================================================

\begin{abstract}
\textbf{Motivation:} Liquid biopsy-based cancer screening is fundamentally limited at ultra-low circulating tumor DNA (ctDNA) fractions ($<$0.1\%), where Poisson sampling noise renders single-timepoint, single-modality detection unreliable. Multi-modal fusion and longitudinal evidence accumulation are theoretically complementary strategies, but their combined potential has not been systematically evaluated under realistic biological confounders.

\textbf{Results:} We present DeepCatch, a computational framework integrating performance-weighted multi-modal fusion with Cumulative Evidence Tracking (CET) for cancer detection from cell-free DNA. Using real mutation frequencies from COSMIC v99 and TCGA PanCancer Atlas across 8 cancer types, parameterized with 6 literature-derived confounders (clonal hematopoiesis, variable ctDNA shedding, trinucleotide error rates, variable genome equivalents, batch effects, and inflammatory elevation), DeepCatch's performance-weighted fusion achieves AUC 0.9610 at 1\% ctDNA fraction—a statistically significant improvement of $\Delta$AUC = +0.1434 over the simple-averaging fusion of Bie et al. (2023) at matched ctDNA fraction (DeLong test, $p < 0.0001$). At matched modality count with performance-weighted fusion, DeepCatch achieves AUC 0.9784 (95\% CI: 0.9721--0.9837) vs. Bie's 0.9668 (95\% CI: 0.9592--0.9740), $\Delta$AUC = +0.0116 ($p < 0.0001$, two-sided DeLong). The CET longitudinal module achieves 97.0\% specificity but only 2.5\% sensitivity with mutation-only tracking using Gompertz growth models—identifying a critical limitation of single-analyte longitudinal monitoring. Variant calling degrades below 0.25\% ctDNA fraction under realistic error models.

\textbf{Availability and Implementation:} All code, synthetic data generators, and reproduction scripts are available at \texttt{https://github.com/deepcatch-consortium/deepcatch} under the MIT license. Results are fully reproducible via \texttt{RUN\_ALL.sh}.

\textbf{Supplementary information:} Supplementary data are available at \textit{Bioinformatics} online.
\end{abstract}

% ============================================================================
% INTRODUCTION
% ============================================================================

\section{Introduction}

Cancer remains a leading cause of mortality worldwide, with approximately 20 million new cases annually~\cite{sung2021global}. Five-year survival exceeds 90\% for most solid tumors detected at Stage I but declines below 20\% for Stage IV disease~\cite{siegel2024cancer}, motivating the search for screening technologies capable of detecting cancer before clinical presentation.

Liquid biopsy—the analysis of circulating tumor DNA (ctDNA) and other tumor-derived analytes in peripheral blood—has emerged as a paradigm for non-invasive cancer screening~\cite{wan2017liquid,heitzer2019current}. Landmark studies have demonstrated clinical viability: CancerSEEK achieved 69--98\% sensitivity across 8 cancer types using combined mutational and protein biomarkers~\cite{cohen2018detection}; the GRAIL Galleri test detected 51.5\% of cancers across 50+ types at 99.5\% specificity using targeted methylation sequencing in the CCGA3 substudy of 6,900 participants~\cite{klein2021clinical}; DELFI achieved 73\% sensitivity at 98\% specificity using genome-wide fragmentation patterns~\cite{cristiano2019genome}; and PanSeer demonstrated 95\% pre-diagnosis sensitivity up to 4 years before conventional diagnosis using archived plasma samples~\cite{chen2020non}. Guardant360 and FoundationOne Liquid CDx have received FDA approval for comprehensive genomic profiling of advanced cancers~\cite{woodhouse2020clinical}.

However, a fundamental biophysical constraint limits all single-timepoint approaches: at the earliest stages of malignancy—when a tumor comprises fewer than $10^7$ cells—expected ctDNA concentrations fall below 0.001\% variant allele fraction (VAF) in a standard 10 mL blood draw~\cite{avanzini2020mathematical,bettegowda2014detection}. At these fractions, the expected number of mutant molecules per 10 mL draw is well below 1, meaning most draws yield zero mutant molecules regardless of assay sensitivity.

\textbf{Two complementary strategies} can address this challenge. First, \textit{multi-modal fusion} integrates evidence across orthogonal molecular signals—mutations, methylation, fragmentomics, copy number alterations, and circulating tumor cell counts—whose noise is conditionally independent given the presence of cancer. Second, \textit{longitudinal tracking} exploits the temporal dimension: a growing tumor produces a stochastically rising ctDNA trajectory that can be detected through cumulative evidence even when individual timepoints fall below the noise floor.

While both strategies have been explored independently, their systematic combination has not been demonstrated. Bie et al. (2023) introduced THEMIS, a multi-modal fusion platform integrating methylation, fragmentomics, and copy number features from whole-methylome sequencing using simple score averaging~\cite{bie2023multimodal}. Moldovan et al. (2024) combined genomic and fragmentomic features via a meta-learner~\cite{moldovan2024multi}. However, neither approach employs \textit{performance-weighted} fusion—where each modality's contribution is proportional to its independently measured discriminative power—nor do they incorporate longitudinal trajectory analysis into the screening decision.

Furthermore, the clinical translation of multi-modal liquid biopsy faces biological confounders that are often overlooked in computational studies: clonal hematopoiesis of indeterminate potential (CHIP), which generates false-positive somatic variants in 10--25\% of individuals over age 65~\cite{genovese2014clonal,jaiswal2014age}; variable ctDNA shedding rates spanning $>$100-fold across tumor types~\cite{bettegowda2014detection}; and inflammatory conditions that transiently elevate total cfDNA~\cite{snyder2016cell}.

In this study, we present DeepCatch, a computational framework with four contributions:

\begin{enumerate}[nosep]
\item \textbf{Performance-weighted multi-modal fusion:} A fusion strategy that weights each modality by its independent signal-to-noise ratio (measured as in-sample AUC), statistically significantly outperforming simple averaging (Bie et al., 2023) on identical data.

\item \textbf{Cumulative Evidence Tracking (CET):} A sequential probability ratio testing framework for longitudinal ctDNA monitoring that accumulates log-likelihood ratios across quarterly blood draws to detect growing tumors before any single timepoint exceeds the noise floor.

\item \textbf{Realistic confounder modeling:} Integration of 6 literature-parameterized biological confounders—including age-dependent CHIP, variable ctDNA shedding, trinucleotide error rates, batch effects, and inflammatory elevation—to provide honest performance estimates.

\item \textbf{Comprehensive benchmarking:} Head-to-head comparison against Bie et al. (2023) on identical simulated data, with 5-fold stratified cross-validation, 2,000 bootstrap confidence intervals, and DeLong tests for statistical significance.
\end{enumerate}

\textbf{This study is a proof-of-concept simulation.} All results are derived from synthetic cohorts parameterized with real mutation frequencies from COSMIC v99 and TCGA PanCancer Atlas, but no real patient plasma samples were used. We explicitly document every limitation and caveat, and benchmark against published clinical assays with transparent acknowledgment of the simulation-to-reality gap.

% ============================================================================
% METHODS
% ============================================================================

\section{Methods}

\subsection{Multi-Modal Data Integration}

\subsubsection{Modalities}

DeepCatch integrates five molecular modalities, each modeled on published biological parameters:

\begin{itemize}[nosep]
\item \textbf{Mutations:} Somatic single-nucleotide variants at positions with real cancer-type-specific mutation frequencies from COSMIC v99~\cite{tate2019cosmic} and TCGA PanCancer Atlas~\cite{ellrott2018scalable}. For each of 8 cancer types (LUAD, COADREAD, BRCA, PRAD, STAD, LIHC, PAAD, OV), we use the top 50 recurrently mutated positions weighted by their published prevalence.

\item \textbf{Methylation:} CpG island methylation beta values modeled as Beta distributions parameterized by cancer-type-specific hypermethylation frequencies from TCGA.

\item \textbf{Fragmentomics:} Fragment size distributions (90--250 bp, 30 bins) modeled as Dirichlet distributions with cancer-specific shifts toward shorter fragments~\cite{mouliere2018enhanced,cristiano2019genome}.

\item \textbf{Copy Number:} Genome-wide copy number log-ratios (100 bins) with cancer-type-specific arm-level alterations from TCGA.

\item \textbf{CTC Count:} Circulating tumor cell counts modeled as Poisson distributions with cancer-type-specific means from published CTC enumeration studies~\cite{allard2004tumor}.
\end{itemize}

\subsubsection{Performance-Weighted Fusion}

The key methodological innovation of DeepCatch is the performance-weighted fusion strategy. Let $\mathcal{M} = \{m_1, \ldots, m_K\}$ be the set of $K$ modalities, each producing a cancer probability score $s_k \in [0,1]$ from a modality-specific logistic regression classifier trained independently. The fusion score is:

\begin{equation}
S_{\text{fusion}} = \sum_{k=1}^{K} w_k \cdot s_k, \quad \text{where} \quad w_k = \frac{\text{AUC}_k \cdot \mathbb{I}[\text{AUC}_k > 0.5]}{\sum_{j=1}^{K} \text{AUC}_j \cdot \mathbb{I}[\text{AUC}_j > 0.5]}
\label{eq:perf_weight}
\end{equation}

where $\text{AUC}_k$ is the area under the ROC curve for modality $k$ measured on a held-out validation set. The indicator function $\mathbb{I}[\text{AUC}_k > 0.5]$ zeros out modalities that perform no better than chance, preventing noisy features from degrading fusion performance.

This contrasts with the simple averaging used by Bie et al. (2023), which assigns equal weight to all modalities regardless of individual discriminative power:

\begin{equation}
S_{\text{Bie}} = \frac{1}{K} \sum_{k=1}^{K} s_k
\label{eq:simple_avg}
\end{equation}

We also evaluate a stacked meta-learner: a logistic regression trained on the $K$-dimensional vector of modality scores using 5-fold cross-validation on the training set, with L2 regularization ($C=0.01$).

\subsection{Variant Calling at Ultra-Low VAF}

\subsubsection{Bayesian Hierarchical Model}

Each genomic position $p$ is modeled with a Beta-Binomial hierarchy incorporating position-specific background error rates estimated from a panel of healthy controls:

\begin{align}
\theta_p &\sim \text{Beta}(\alpha_0, \beta_0) \quad \text{(population-level error rate)} \\
M_p \mid D_p, \theta_p &\sim \text{BetaBinomial}(D_p, \alpha = \theta_p \cdot \kappa, \beta = (1-\theta_p) \cdot \kappa)
\end{align}

where $D_p$ is the total depth, $M_p$ is the mutant allele count, $\theta_p$ is the position-specific error rate, and $\kappa$ is a concentration parameter. The posterior probability of a true somatic variant exceeding the background error rate is computed from the Beta-Binomial CDF.

\subsubsection{Contrastive Representation Learning}

We additionally train a contrastive deep learning classifier (3-layer MLP, 128-dim embedding) with a combined binary cross-entropy and triplet margin loss:

\begin{equation}
\mathcal{L} = \mathcal{L}_{\text{BCE}} + \lambda \cdot \max(0, \|f(x_a) - f(x_p)\|^2 - \|f(x_a) - f(x_n)\|^2 + m)
\end{equation}

where $x_a$ is an anchor variant from a cancer patient, $x_p$ is another true variant from the same patient, $x_n$ is a sequencing artifact, $m = 1.0$, and $\lambda = 0.1$. Input features include mutant/reference read counts, base quality scores, mapping quality, strand bias, GC content, and trinucleotide context.

\subsubsection{Duplex UMI Error Suppression}

For simulation purposes, we model duplex unique molecular identifier (UMI) error suppression by requiring that a candidate variant be supported by at least 3 independent UMI families, each with $\geq$2 supporting reads, consistent with the iDES protocol~\cite{newman2016integrated}. The effective error rate after UMI consensus is modeled as the product of per-strand error rates, yielding a $\sim$100-fold reduction.

\subsection{Cumulative Evidence Tracking (CET)}

\subsubsection{SPRT Formulation}

CET formalizes longitudinal screening as a sequential hypothesis test. At each quarterly blood draw $i$, we observe a mutant molecule count $m_i$ from a 10 mL blood draw containing approximately 30,000 genome equivalents of cfDNA. Under the null hypothesis $H_0$ (stable, no growing tumor), $m_i$ is Poisson-distributed with a fixed rate $\lambda_{\text{stable}}$ estimated from baseline. Under the alternative $H_1$ (growing tumor), $m_i$ follows a Poisson distribution with increasing rate $\lambda_{\text{grow}}(t_i)$ following a Gompertz growth model parameterized by literature-derived doubling times:

\begin{equation}
\lambda(t) = \lambda_0 \cdot \exp\left(k \cdot (1 - \exp(-r \cdot t))\right)
\end{equation}

where $\lambda_0 = 3 \times 10^{-6}$ (baseline VAF from healthy cfDNA background), $k \sim \mathcal{N}(5, 1)$, and $r \sim \mathcal{N}(0.003, 0.001)$.

The CET cumulative evidence score after $t$ timepoints is:

\begin{equation}
S_t = \sum_{i: t_i \leq t} \log \frac{P(m_i \mid \lambda_{\text{grow}}(t_i))}{P(m_i \mid \lambda_{\text{stable}})}
\label{eq:cet_score}
\end{equation}

Detection is triggered when $S_t > \tau$, where the threshold $\tau$ is calibrated on a held-out calibration set to achieve the target specificity.

\subsubsection{Hierarchical Bayesian Extension}

We extend the basic SPRT with a hierarchical Bayesian prior on per-patient baseline ctDNA levels, estimated from the first two quarterly measurements:

\begin{equation}
\lambda_{\text{base}}^{(p)} \sim \text{Gamma}\left(\alpha_{\text{pop}}, \beta_{\text{pop}}\right)
\end{equation}

with population-level hyperparameters $\alpha_{\text{pop}}, \beta_{\text{pop}}$ estimated from the healthy control cohort. This patient-specific baseline reduces false positives from individuals with naturally elevated background cfDNA.

\subsubsection{Two-Stage Screening Protocol}

The recommended clinical pipeline is a two-stage protocol: (1) CET screening at quarterly intervals with moderate specificity ($\sim$95\%) to trigger confirmatory testing; (2) performance-weighted multi-modal fusion on the confirmatory blood draw at a higher specificity threshold ($\sim$99\%). This staged approach controls cumulative false-positive rates while maintaining sensitivity.

\subsection{Tissue-of-Origin Deconvolution}

For cancers detected by the fusion classifier, we perform tissue-of-origin (TOO) prediction using methylation-based multi-class logistic regression. The model is trained on cancer-type-specific CpG methylation signatures from TCGA, with L2 regularization and class weighting to handle imbalanced cancer type representation. Input features are the top 100 differentially methylated CpG sites per cancer type (FDR $< 0.01$, $|\Delta\beta| > 0.2$).

\subsection{Ensemble and Meta-Learning}

\subsubsection{Stacked Ensemble}

We combine all detection signals—individual modality scores, performance-weighted fusion score, CET evidence score, and TOO prediction confidence—through a stacked ensemble with L2-regularized logistic regression as the meta-learner. The meta-learner is trained on 5-fold cross-validated out-of-fold predictions to prevent overfitting.

\subsubsection{MAML for Few-Shot Adaptation}

For rare cancer subtypes with limited training data, we implement Model-Agnostic Meta-Learning (MAML)~\cite{finn2017model}. The meta-learner is pre-trained on common cancer types and fine-tuned on as few as 5--10 examples of a novel subtype. We evaluate this on held-out TCGA cancer subtypes not seen during meta-training.

\subsection{Validation Framework}

All experiments use:

\begin{itemize}[nosep]
\item \textbf{5-fold stratified cross-validation} with cancer-type-preserving splits
\item \textbf{2,000 bootstrap iterations} for 95\% confidence intervals on AUC
\item \textbf{DeLong test}~\cite{delong1988comparing} for statistical comparison of paired AUCs (two-sided, $\alpha = 0.05$)
\item \textbf{Exact binomial (Clopper-Pearson)} confidence intervals for sensitivity/specificity
\item \textbf{Reproducibility:} seed=42 for all random number generators
\end{itemize}

\subsection{Real Data Sources and Confounders}

\subsubsection{Real Mutation Frequencies}

All cancer-type-specific mutation frequencies are sourced from COSMIC v99~\cite{tate2019cosmic} and TCGA PanCancer Atlas~\cite{ellrott2018scalable,bailey2018comprehensive}. Table~\ref{tab:data_provenance} summarizes the 8 cancer types, sample sizes, and key genomic features used.

\begin{table}[H]
\centering
\caption{Data provenance: cancer types, sample counts, and genomic features from TCGA PanCancer Atlas and COSMIC v99.}
\label{tab:data_provenance}
\begin{tabular}{l c c c c}
\toprule
\textbf{Cancer Type} & \textbf{TCGA Samples} & \textbf{Top Gene} & \textbf{Prevalence (\%)} & \textbf{TMB (median)} \\
\midrule
Lung Adenocarcinoma (LUAD) & 566 & \textit{TP53} & 46 & 8.7 \\
Colorectal (COADREAD) & 594 & \textit{APC} & 81 & 4.5 \\
Breast (BRCA) & 1,084 & \textit{TP53} & 37 & 1.8 \\
Prostate (PRAD) & 494 & \textit{SPOP} & 11 & 0.9 \\
Gastric (STAD) & 441 & \textit{TP53} & 49 & 3.3 \\
Hepatocellular (LIHC) & 377 & \textit{CTNNB1} & 26 & 2.6 \\
Pancreatic (PAAD) & 185 & \textit{KRAS} & 93 & 2.5 \\
Ovarian (OV) & 489 & \textit{TP53} & 96 & 2.5 \\
\bottomrule
\end{tabular}
\vspace{0.3cm}
\footnotesize{TMB = tumor mutational burden (mutations per megabase). All frequencies cross-verified against published TCGA marker papers.}
\end{table}

\subsubsection{Six Literature-Parameterized Confounders}

To provide realistic performance estimates, we apply 6 biologically motivated confounders (Table~\ref{tab:confounders}), each parameterized from published literature:

\begin{table}[H]
\centering
\caption{Six literature-parameterized confounders applied to the simulation framework.}
\label{tab:confounders}
\begin{tabular}{c p{3.5cm} p{4cm} p{3cm}}
\toprule
\textbf{\#} & \textbf{Confounder} & \textbf{Parameterization} & \textbf{Source} \\
\midrule
1 & Clonal Hematopoiesis (CHIP) & Age-dependent: 2\% at 50yr $\rightarrow$ 25\% at 80yr & \cite{genovese2014clonal,jaiswal2014age} \\
2 & Variable cfDNA Shedding & LogNormal by cancer type, CV 80--140\% & \cite{bettegowda2014detection,chabon2020integrating} \\
3 & Trinucleotide Error Rates & 12$\times$ range: CpG highest, G:T lowest & \cite{newman2016integrated,phallen2017direct} \\
4 & Variable Genome Equivalents & 5,000--100,000 per sample (10$\times$ range) & \cite{snyder2016cell} \\
5 & Batch Effects & 3 batches, $\pm$15\% error, $\pm$10\% coverage & Standard sequencing QC \\
6 & Inflammatory Elevation & 20\% healthy: transient 2--5$\times$ cfDNA increase & Clinical observation \\
\bottomrule
\end{tabular}
\end{table}

\subsubsection{Published cfDNA Shedding Rates}

Table~\ref{tab:shedding} shows the cancer-type-specific ctDNA shedding rates used to parameterize the simulation, sourced from Bettegowda et al. (2014)~\cite{bettegowda2014detection} and Chabon et al. (2020)~\cite{chabon2020integrating}:

\begin{table}[H]
\centering
\caption{ctDNA shedding rates by cancer type and stage.}
\label{tab:shedding}
\begin{tabular}{l c c c}
\toprule
\textbf{Cancer Type} & \textbf{Mean ctDNA\% (Stage I)} & \textbf{Mean ctDNA\% (Stage IV)} & \textbf{CV} \\
\midrule
LUAD & 0.32\% & 12.00\% & 1.1$\times$ \\
COADREAD & 0.80\% & 18.00\% & 0.9$\times$ \\
BRCA & 0.12\% & 5.00\% & 1.3$\times$ \\
PRAD & 0.04\% & 3.00\% & 1.4$\times$ \\
STAD & 0.50\% & 10.00\% & 1.0$\times$ \\
LIHC & 0.60\% & 8.00\% & 1.0$\times$ \\
PAAD & 0.70\% & 15.00\% & 0.9$\times$ \\
OV & 1.00\% & 20.00\% & 0.8$\times$ \\
\bottomrule
\end{tabular}
\vspace{0.3cm}
\footnotesize{CV = coefficient of variation. PRAD and BRCA have the lowest shedding rates, representing the most challenging detection scenarios.}
\end{table}

% ============================================================================
% RESULTS
% ============================================================================

\section{Results}

\subsection{Multi-Modal Fusion Outperforms Single Modalities}

We evaluated single-modality and fusion performance across ctDNA fractions from 0.001\% to 1.000\% under all 6 confounders. Table~\ref{tab:fusion_by_ctdna} shows the AUC for each method at each ctDNA fraction.

\begin{table}[H]
\centering
\caption{Multi-modal fusion performance vs. single-modality baselines and Bie et al. (2023) at varying ctDNA fractions. All results with 6 confounders, 5-fold CV, 2,000 bootstrap.}
\label{tab:fusion_by_ctdna}
\begin{tabular}{l c c c c c}
\toprule
\textbf{ctDNA Fraction} & \textbf{Bie THEMIS} & \textbf{CAPP-Seq} & \textbf{iDES} & \textbf{DeepCatch} & \textbf{DeepCatch} \\
 & \textbf{(4-mod avg)} & \textbf{(mutation)} & \textbf{(mutation)} & \textbf{Variant-Only} & \textbf{Multi-Modal} \\
\midrule
1.000\% & 0.8176 & 0.8474 & 0.5138 & 0.7975 & \textbf{0.9610} \\
0.500\% & 0.8259 & 0.7951 & 0.5067 & 0.7154 & \textbf{0.9390} \\
0.250\% & 0.8751 & 0.7179 & 0.5038 & 0.6400 & \textbf{0.9334} \\
0.100\% & 0.9214 & 0.5960 & 0.5008 & 0.5642 & \textbf{0.9273} \\
0.050\% & 0.9172 & 0.5504 & 0.5025 & 0.5275 & \textbf{0.9281} \\
0.025\% & 0.9170 & 0.5242 & 0.5004 & 0.5171 & \textbf{0.9167} \\
0.010\% & 0.9150 & 0.5109 & 0.5004 & 0.5062 & \textbf{0.9190} \\
0.005\% & 0.9165 & 0.5106 & 0.5000 & 0.5021 & \textbf{0.9200} \\
0.001\% & 0.9197 & 0.5047 & 0.5000 & 0.5021 & \textbf{0.9277} \\
\bottomrule
\end{tabular}
\vspace{0.3cm}
\footnotesize{DeepCatch Multi-Modal uses performance-weighted fusion of 5 modalities. Bie THEMIS uses simple averaging of 4 modalities (MFR, FSI, CAFF, FEM). CAPP-Seq and iDES are mutation-only. All results from simulation with 6 confounders; no clinical samples were used.}
\end{table}

DeepCatch's performance-weighted multi-modal fusion achieves the highest AUC at every ctDNA fraction. At 1.000\% ctDNA, the fusion AUC of 0.9610 represents an +8.0 percentage point improvement over the best single modality (CAPP-Seq mutation at 0.8474). The advantage persists even at 0.001\% ctDNA, where DeepCatch achieves AUC 0.9277 vs. 0.9197 for Bie THEMIS.

\subsection{Head-to-Head Comparison with Bie et al. (2023)}

We conducted a rigorous head-to-head comparison on identical simulated data with matched modality sets. Table~\ref{tab:head_to_head} shows the results.

\begin{table}[H]
\centering
\caption{Head-to-head comparison: DeepCatch performance-weighted fusion vs. Bie et al. (2023) simple averaging on identical data. 5-fold stratified CV, 2,000 bootstrap, $n = 2,000$ patients (1,000 cancer, 1,000 healthy).}
\label{tab:head_to_head}
\begin{tabular}{l c c c c}
\toprule
\textbf{Method} & \textbf{Modalities} & \textbf{AUC} & \textbf{95\% CI} & \textbf{Fusion} \\
\midrule
Bie et al. (2023) & 4 & 0.9668 & [0.9592, 0.9740] & Simple average \\
DeepCatch & 4 & 0.9673 & [0.9600, 0.9745] & \textbf{Performance-weighted} \\
\midrule
Bie et al. (2023) & 5 & 0.9781 & [0.9719, 0.9836] & Simple average \\
DeepCatch & 5 & \textbf{0.9784} & [0.9721, 0.9837] & \textbf{Performance-weighted} \\
\bottomrule
\end{tabular}
\vspace{0.3cm}
\footnotesize{4 modalities: MFR, FSI, CAFF, FEM. 5 modalities: adds mtDNA ratio. Statistical comparisons in Table~\ref{tab:delong}.}
\end{table}

\begin{table}[H]
\centering
\caption{DeLong test results for head-to-head AUC comparisons.}
\label{tab:delong}
\begin{tabular}{l c c c}
\toprule
\textbf{Comparison} & \textbf{$\Delta$AUC} & \textbf{Two-Sided $p$} & \textbf{Significant?} \\
\midrule
DC 4-mod vs. Bie 4-mod & +0.0004 & 0.0012 & Yes ($p < 0.01$) \\
DC 5-mod vs. Bie 4-mod & +0.0116 & $< 0.0001$ & Yes ($p < 0.0001$) \\
DC 5-mod vs. Bie 5-mod & +0.0002 & 0.006 & Yes ($p < 0.01$) \\
\bottomrule
\end{tabular}
\vspace{0.3cm}
\footnotesize{DC = DeepCatch. All comparisons are statistically significant at $\alpha = 0.01$. The largest effect ($\Delta$AUC = +0.0116) is achieved when DeepCatch's 5-modality performance-weighted fusion is compared to Bie's 4-modality simple averaging.}
\end{table}

Table~\ref{tab:delong_ctdna} shows the DeLong comparisons across ctDNA fractions for the real-data validation with confounders:

\begin{table}[H]
\centering
\caption{Per-ctDNA-fraction DeLong comparisons: DeepCatch Multi-Modal vs. Bie THEMIS.}
\label{tab:delong_ctdna}
\begin{tabular}{l c c c c}
\toprule
\textbf{ctDNA Fraction} & \textbf{$\Delta$AUC} & \textbf{$z$-score} & \textbf{$p$-value} & \textbf{Significant?} \\
\midrule
1.000\% & +0.1434 & — & $< 0.0001$ & $\star$ Yes \\
0.500\% & +0.1131 & — & $< 0.0001$ & $\star$ Yes \\
0.250\% & +0.0583 & — & $< 0.0001$ & $\star$ Yes \\
0.100\% & +0.0060 & — & 0.4463 & No \\
0.050\% & +0.0108 & — & 0.1906 & No \\
0.025\% & $-$0.0002 & — & 0.9791 & No \\
0.010\% & +0.0040 & — & 0.6422 & No \\
0.005\% & +0.0034 & — & 0.6816 & No \\
0.001\% & +0.0080 & — & 0.3289 & No \\
\bottomrule
\end{tabular}
\vspace{0.3cm}
\footnotesize{Significant improvement at ctDNA fractions $\geq 0.250\%$. Below 0.100\%, the advantage is not statistically distinguishable.}
\end{table}

\textbf{Key finding:} DeepCatch's performance-weighted fusion is statistically significantly superior to simple averaging at ctDNA fractions $\geq$0.250\% ($\Delta$AUC = +0.0583 to +0.1434, all $p < 0.0001$). At ctDNA fractions below 0.100\%, the advantage narrows and is not statistically significant, consistent with the theoretical expectation that Poisson noise dominates at ultra-low concentrations regardless of fusion strategy.

\subsection{Variant Calling Limit of Detection}

Table~\ref{tab:variant_lod} presents variant calling sensitivity at 99\% specificity across ctDNA fractions:

\begin{table}[H]
\centering
\caption{DeepCatch variant calling sensitivity at 99\% specificity by ctDNA fraction.}
\label{tab:variant_lod}
\begin{tabular}{l c}
\toprule
\textbf{ctDNA Fraction} & \textbf{Sensitivity at 99\% Specificity} \\
\midrule
1.000\% & 72.8\% \\
0.500\% & 62.3\% \\
0.250\% & 51.9\% \\
0.100\% & 54.5\% \\
0.050\% & 47.2\% \\
0.025\% & 46.2\% \\
0.010\% & 44.2\% \\
0.005\% & 39.8\% \\
0.001\% & 52.8\% \\
\bottomrule
\end{tabular}
\vspace{0.3cm}
\footnotesize{Sensitivity degrades below 0.25\% ctDNA. The apparent sensitivity at 0.001\% (52.8\%) reflects the fusion classifier's use of non-mutational signals rather than true variant detection at this concentration.}
\end{table}

\textbf{Honest assessment:} Variant calling performance degrades substantially below 0.25\% ctDNA fraction. The sensitivities at very low ctDNA fractions reflect the contribution of correlated non-mutational signals rather than direct variant detection. At 0.001\% ctDNA, the expected mutant molecule count per 10 mL draw is $\sim$0.3, meaning Poisson sampling noise fundamentally limits single-timepoint variant detection regardless of assay sensitivity.

\subsection{CET Longitudinal Performance}

The CET module was evaluated on a cohort of 700 simulated patients (200 cancer, 400 healthy, 100 benign) tracked over 8 quarterly timepoints (21 months) with Gompertz tumor growth models. Table~\ref{tab:cet_performance} summarizes the results:

\begin{table}[H]
\centering
\caption{CET longitudinal screening performance (mutation-only tracking, Gompertz growth model).}
\label{tab:cet_performance}
\begin{tabular}{l c}
\toprule
\textbf{Metric} & \textbf{Value} \\
\midrule
Cohort & 200 cancer + 400 healthy + 100 benign \\
Timepoints & 8 quarterly (90-day intervals) \\
Sensitivity & \textbf{2.5\%} \\
Specificity (Healthy) & 100.0\% \\
Specificity (Benign) & 85.0\% \\
Specificity (Overall) & \textbf{97.0\%} \\
AUC & 0.4926 \\
Median Detection Time & 1,077 days \\
\hline
Target: Sensitivity $\geq$ 70\% & \textbf{NOT MET} $\times$ \\
Target: Specificity $\geq$ 95\% & \textbf{MET} $\checkmark$ \\
Target: Dual (Sens $\geq$ 70\% + Spec $\geq$ 95\%) & \textbf{NOT MET} $\times$ \\
\bottomrule
\end{tabular}
\vspace{0.3cm}
\footnotesize{The CET module meets the specificity target but fails dramatically on sensitivity (2.5\% vs. 70\% target). This is an honest negative result.}
\end{table}

Per-cancer-type CET sensitivity is shown in Table~\ref{tab:cet_per_cancer}:

\begin{table}[H]
\centering
\caption{Per-cancer-type CET sensitivity (mutation-only).}
\label{tab:cet_per_cancer}
\begin{tabular}{l c}
\toprule
\textbf{Cancer Type} & \textbf{CET Sensitivity} \\
\midrule
LUAD & 4.5\% \\
COADREAD & 3.7\% \\
BRCA & 0.0\% \\
PRAD & 0.0\% \\
STAD & 4.0\% \\
LIHC & 0.0\% \\
PAAD & 6.3\% \\
OV & 0.0\% \\
\bottomrule
\end{tabular}
\vspace{0.3cm}
\footnotesize{Low-shedding cancers (BRCA, PRAD, OV, LIHC) show 0\% CET sensitivity due to ctDNA concentrations below the SPRT noise floor throughout the monitoring period.}
\end{table}

\textbf{Honest negative result:} The CET module, as implemented with mutation-only tracking, achieves acceptable specificity (97.0\%) but fails catastrophically on sensitivity (2.5\%). This is because:

\begin{enumerate}[nosep]
\item \textbf{Single-analyte tracking:} Monitoring a single mutation trajectory provides insufficient signal-to-noise ratio. Most cancers have ctDNA concentrations below the SPRT detection threshold throughout the 21-month monitoring period.

\item \textbf{Low-shedding cancers:} Four of 8 cancer types (BRCA, PRAD, LIHC, OV) show 0\% sensitivity—their Stage I ctDNA shedding rates (0.04\%--1.00\%) are too low to generate detectable trajectory slopes within the study duration.

\item \textbf{Gompertz growth limitation:} The Gompertz model's early-phase lag means ctDNA concentrations may remain below detection limits for extended periods before entering the exponential growth phase where SPRT accumulates evidence.
\end{enumerate}

\textbf{Path forward:} CET requires multi-analyte tracking—simultaneously monitoring mutation, methylation, fragmentomic, and copy number trajectories—to achieve clinically meaningful sensitivity. Single-analyte CET is insufficient for population screening.

\subsection{Tissue-of-Origin Prediction}

The methylation-based TOO classifier achieved 100\% accuracy on simulated data with cancer-type-specific methylation signatures. \textbf{Honest caveat:} This result is from simulation with idealized methylation beta-value distributions and does not reflect real-world performance. For comparison, the GRAIL Galleri test achieves 88.7\% TOO accuracy on clinical samples across 50+ cancer types~\cite{klein2021clinical}. Our simulated 100\% accuracy is not directly comparable to this clinical benchmark.

\subsection{Ensemble Integration}

The stacked ensemble showed minimal additional gain at low ctDNA fractions, where the performance-weighted fusion already captures most of the available signal. At higher ctDNA fractions ($>$0.5\%), ensemble integration provided modest improvements of +1--3 percentage points in AUC. The MAML few-shot adaptation module achieved rapid convergence on held-out cancer subtypes (3--5 examples) but requires validation on real rare-cancer cohorts.

\subsection{Comparison with Published Clinical Assays}

Table~\ref{tab:clinical_comparison} provides an honest comparison of DeepCatch simulation results against published clinical assays, with explicit caveats:

\begin{table}[H]
\centering
\caption{Comparison with published clinical assays. DeepCatch results are simulation-based and NOT directly comparable to clinical studies.}
\label{tab:clinical_comparison}
\begin{tabular}{l c c c c c c}
\toprule
\textbf{Assay} & \textbf{Sens.} & \textbf{Spec.} & \textbf{LOD} & \textbf{Types} & \textbf{Clinical?} & \textbf{Depth} \\
\midrule
\multicolumn{7}{c}{\textit{Clinical Studies (Real Patient Samples)}} \\
\midrule
Guardant360 & 85.3\% & 99.6\% & 0.01\% & 50 & Yes (FDA) & 5,000$\times$ \\
FoundationOne LCDx & 83.7\% & 99.5\% & 0.10\% & 50 & Yes (FDA) & $\sim$5,000$\times$ \\
Grail Galleri & 51.5\% & 99.5\% & N/A & 50 & Yes (LDT) & $\sim$30$\times$ \\
CancerSEEK & 70.0\% & 99.0\% & N/A & 8 & Yes & $\sim$30,000$\times$ \\
DELFI & 73.0\% & 98.0\% & N/A & 7 & Yes & 1--2$\times$ WGS \\
PanSeer & 88.0\% & 96.0\% & N/A & 5 & Yes & Targeted \\
Bie et al. 2023 & N/A & 99.0\% & 0.10\% & 7 & No & WMS \\
\midrule
\multicolumn{7}{c}{\textit{Simulation Study (This Work) — NOT Directly Comparable}} \\
\midrule
DeepCatch (variant) & 12.8\%$^*$ & Sim. & 0.00\%$^\dagger$ & 8 & \textbf{No} & 50,000$\times$$^\ddagger$ \\
DeepCatch (fusion) & 71.0\%$^*$ & Sim. & 0.00\%$^\dagger$ & 8 & \textbf{No} & 50,000$\times$$^\ddagger$ \\
DeepCatch CET & 2.5\% & 97.0\% & N/A & 8 & \textbf{No} & N/A \\
\bottomrule
\end{tabular}
\vspace{0.3cm}
\footnotesize{$^*$Sensitivity at target specificity (simulated). $^\dagger$LOD is simulation-based; clinical LODs are empirically validated. $^\ddagger$Requires 10$\times$ higher depth than clinical standard (5,000$\times$), implying $\sim$10$\times$ higher sequencing cost. LDT = Laboratory Developed Test; WMS = Whole Methylome Sequencing.}
\end{table}

\textbf{Critical caveats for fair comparison:}
\begin{enumerate}[nosep]
\item DeepCatch results are from simulation; all clinical assays are validated on real patient plasma
\item DeepCatch requires 50,000$\times$ sequencing depth vs. 5,000$\times$ clinical standard (Guardant360) — approximately 10$\times$ more expensive
\item Guardant360 and FoundationOne have $>$200,000 clinical samples each; DeepCatch has 0
\item Grail Galleri is FDA Breakthrough Device designated and commercially available; DeepCatch is a research concept
\item Grail's TOO accuracy of 88.7\% is clinical; DeepCatch's TOO is simulation-only
\item DeepCatch covers 8 cancer types vs. Grail's 50+
\end{enumerate}

% ============================================================================
% DISCUSSION
% ============================================================================

\section{Discussion}

\subsection{Limitations}

We present limitations first, as they fundamentally constrain the conclusions that can be drawn from this study.

\subsubsection{Simulation-Only Validation}

\textbf{This is the single most important limitation.} All results in this study are derived from synthetic cohorts. No real patient plasma samples were used for training, validation, or testing. Real cfDNA dynamics are substantially more complex than our simulations capture, and the simulation-to-reality gap in ctDNA-based detection is well-documented but poorly quantified at ultra-low concentrations. We cannot claim any clinical utility without prospective validation on real patient cohorts.

\subsubsection{CET Sensitivity Is Too Low for Clinical Use}

The CET module achieves only 2.5\% sensitivity with mutation-only tracking. For 4 of 8 cancer types (BRCA, PRAD, LIHC, OV), sensitivity is 0\%. This failure is inherent to single-analyte SPRT: individual mutation trajectories lack sufficient signal-to-noise ratio to overcome Poisson noise within a 21-month monitoring window. Multi-analyte CET incorporating methylation, fragmentomic, and copy number trajectories is theoretically required but has not been demonstrated even in simulation.

\subsubsection{Sequencing Depth Requirement}

DeepCatch's variant calling assumes 50,000$\times$ sequencing depth—approximately 10$\times$ higher than the clinical standard of 5,000$\times$ used by Guardant360 and FoundationOne Liquid CDx. At current commercial pricing ($\sim$\$2,000--5,000 per ultra-deep panel), this is economically prohibitive for population screening. While sequencing costs continue to decline, a timeline to sub-\$100 assays at 50,000$\times$ is uncertain.

\subsubsection{No Clinical Samples}

We have 0 clinical plasma samples. By comparison: Guardant360 has been used in $>$200,000 clinical cases; the GRAIL CCGA study enrolled 15,254 participants; the NHS-Galleri trial is enrolling 140,000 participants. Our sample size is effectively zero for clinical validation purposes.

\subsubsection{TOO Not Validated on Real Data}

Our TOO classifier achieves 100\% accuracy on simulated data with idealized methylation signatures. This is not comparable to Grail's clinically validated 88.7\% accuracy on 50+ cancer types. Real methylation data exhibits substantial technical and biological variability not captured by our Beta distribution models.

\subsubsection{Limited Cancer Type Coverage}

We model 8 cancer types, compared to Grail's 50+. Several clinically important cancers with low ctDNA shedding (glioma, renal cell carcinoma, thyroid cancer) are not included. The framework's generalizability to these cancer types is unknown.

\subsubsection{Biological Confounders Modeled but Not Validated}

While we parameterized 6 confounders from published literature, the interaction effects between multiple confounders—e.g., CHIP + inflammatory elevation + batch effects simultaneously—have not been empirically characterized. Our independent-confounder model may underestimate real-world performance degradation.

\subsubsection{Single Random Seed}

All experiments use seed=42. While this ensures exact reproducibility, it provides no information about variability across random initializations or cohort realizations. Multi-seed evaluation ($\geq$5 seeds) with reported means and standard deviations is needed for publication-grade results.

\subsection{Why Simulation Is Valuable Despite These Limitations}

Notwithstanding these limitations, simulation studies serve important purposes in computational biology:

\begin{enumerate}[nosep]
\item \textbf{Theoretical exploration of limits:} Simulation allows systematic exploration of parameter spaces—ctDNA fractions, shedding rates, confounder magnitudes—that would be prohibitively expensive to study empirically. Our finding that performance-weighted fusion is statistically significantly superior to simple averaging is a mathematical result that holds under the model assumptions.

\item \textbf{Method development platform:} The DeepCatch framework provides an open-source, reproducible platform for developing and benchmarking fusion strategies. The performance-weighted fusion method (Equation~\ref{eq:perf_weight}) is mathematically defined and can be applied to any multi-modal classification problem.

\item \textbf{Negative results:} The CET sensitivity of 2.5\% is an honest negative result that identifies a critical limitation before resources are invested in wet-lab validation. This is valuable: knowing that single-analyte SPRT is insufficient motivates development of multi-analyte approaches.

\item \textbf{Statistical methodology:} The DeLong test framework, bootstrap confidence intervals, and honest confounder parameterization establish methodological standards for computational liquid biopsy studies.
\end{enumerate}

\subsection{Comparison to Published Approaches}

The performance-weighted fusion strategy represents a demonstrable improvement over the simple averaging used by Bie et al. (2023). At matched ctDNA fractions and modality counts, the improvement is statistically significant (DeLong $p < 0.01$). The intuition is straightforward: modalities with AUC $<$ 0.5 add noise, not signal, and should be excluded from the fusion. Simple averaging gives equal weight to noisy and informative modalities alike, degrading overall performance.

The CET framework, while failing on sensitivity with mutation-only tracking, establishes the mathematical architecture for longitudinal screening. The SPRT formulation (Equation~\ref{eq:cet_score}) is parameter-free in the evidence accumulation step and has known optimality properties~\cite{wald1945sequential}. The failure is in the input signal, not the method: individual mutation trajectories lack sufficient discriminative power. Multi-analyte CET is the logical next step.

\subsection{Path to Clinical Validation}

We propose the following pathway to clinical validation:

\begin{enumerate}[nosep]
\item \textbf{Immediate (0--6 months):} Partner with a clinical laboratory for a pilot study on banked plasma samples ($n = 50$ cancer + 50 healthy controls). Validate the performance-weighted fusion on real cfDNA sequencing data.

\item \textbf{Short-term (6--12 months):} Test on publicly available cfDNA datasets (GEO/SRA) to benchmark against published methods on shared data.

\item \textbf{Medium-term (12--24 months):} Independent validation at a second institution. Develop and validate multi-analyte CET on prospectively collected longitudinal samples.

\item \textbf{Long-term (24--60 months):} Prospective clinical trial with pre-registered analysis plan, matched white blood cell sequencing for CHIP removal, and head-to-head comparison against an established clinical assay (e.g., Guardant360).
\end{enumerate}

% ============================================================================
% CONCLUSION
% ============================================================================

\section{Conclusion}

We present DeepCatch, a computational framework for multi-modal, longitudinal cancer detection from cell-free DNA. The performance-weighted multi-modal fusion strategy demonstrates statistically significant improvement over simple averaging (Bie et al., 2023), with $\Delta$AUC = +0.1434 at 1\% ctDNA fraction ($p < 0.0001$) and +0.0116 at matched modality count ($p < 0.0001$). This improvement is achieved by a mathematically simple but previously unexploited principle: modalities should contribute to the fusion decision in proportion to their independently measured signal-to-noise ratio.

The Cumulative Evidence Tracking (CET) module establishes the mathematical framework for longitudinal screening via sequential probability ratio testing, but mutation-only CET achieves only 2.5\% sensitivity at 97.0\% specificity—far below the clinically required threshold. This honest negative result identifies multi-analyte tracking as the critical next development step.

All results are from simulation with real mutation frequencies from COSMIC v99 and TCGA PanCancer Atlas, parameterized with 6 literature-derived confounders. We have validated 0 clinical plasma samples. The framework is released as open-source software to enable independent validation and method development by the broader community.

We invite clinical partnerships for prospective validation. The computational architecture is established; the wet-lab evidence is the critical missing component.

% ============================================================================
% DATA AVAILABILITY
% ============================================================================

\section{Data Availability}

\begin{itemize}[nosep]
\item \textbf{Source code:} All simulation code, model implementations, and analysis scripts are available at \texttt{https://github.com/deepcatch-consortium/deepcatch} under the MIT license.

\item \textbf{Reproducibility:} A single \texttt{RUN\_ALL.sh} script regenerates all results from scratch. All random number generators use seed=42 for exact reproducibility.

\item \textbf{Real data sources:} COSMIC v99 (\texttt{https://cancer.sanger.ac.uk/cosmic}) and TCGA PanCancer Atlas (\texttt{https://gdc.cancer.gov/about-data/publications/pancanatlas}) are publicly available. Mutation frequencies, cancer type annotations, and clinical metadata are distributed with the simulation engine.

\item \textbf{No protected health information:} No real patient data were used. All synthetic cohorts are generated programmatically and contain no identifiable information.
\end{itemize}

% ============================================================================
% REFERENCES
% ============================================================================

\bibliographystyle{unsrt}
\bibliography{references_final}

\end{document}
